\chapter{Materials \& Methods}
\label{ch:methods}

\section{Data Source and Curation}
The datasets utilized in this meta-analysis include three study entries from the NASA Open Science Data Repository (OSDR): OSD-745 (microbiological and nutritional profiles for space-grown lettuce), OSD-655 (nutritional characterization of space-grown Mizuna mustard), and OSD-780 (microbial plate count assays for Mizuna mustard). All files and associated metadata were programmatically fetched from the OSDR API (\url{https://osdr.nasa.gov}) utilizing standard REST queries and parsed from ISA-Tab (Investigation, Study, Assay) format sheets. 

\begin{coseagronomy}[FAIR Principles]
The data harmonization pipeline developed for this project aligns with FAIR principles. All raw assay files, metadata sheets, microbiological plate counts, and biochemical readouts across studies were merged into a consolidated, version-controlled analytical matrix.
\end{coseagronomy}

\section{Plant Material and Growth Conditions}
Two leafy green crops were analyzed: \textit{Lactuca sativa} cv. 'Outredgeous' (red romaine lettuce) grown during experimental epochs VEG-01A (May-Jun 2014), VEG-01B (Jul-Aug 2015), and VEG-03A (Oct-Dec 2016), and \textit{Brassica rapa} cv. 'Tokyo Bekana' (Mizuna mustard) grown during VEG-04A and VEG-04B. Both crops were seeded into Veggie "pillows" filled with an arcillite-based growth substrate and controlled-release fertilizer (Nutricote 18-6-8, type 180).

Cultivation occurred aboard the ISS (Flight) or under KSC ground environment-control chambers (Ground Control). The Veggie LED light cap provided Red (630 nm), Blue (455 nm), and Green (530 nm) lighting, with photoperiods of 16-hour light / 8-hour dark. Telemetry data from the ISS cabin (relative humidity, temperature, and carbon dioxide concentration) was mirrored in the ground control chambers with a 24- to 72-hour delay to isolate microgravity effects.

\section{Sample Collection and Analytical Methods}
Upon reaching maturity (approximately 33 days post-imbibition for lettuce and 28 days for Mizuna), crops were harvested by the crew. A portion of the fresh tissue was consumed in flight, while the remainder was immediately flash-frozen at $-80^\circ$C using the MELFI (Minus Eighty-Degree Laboratory Freezer for ISS) and subsequently returned to Earth aboard SpaceX Dragon capsules for detailed laboratory analyses.

\subsection{Elemental Profiling (ICP-OES)}
Lyophilized leaf tissue was digested in nitric acid and analyzed for macronutrients and micronutrients using Inductively Coupled Plasma Optical Emission Spectrometry (ICP-OES) via a Thermo Scientific iCAP 7000 series spectrometer. Elements quantified included Iron (Fe), Potassium (K), Sodium (Na), Phosphorus (P), Sulfur (S), Zinc (Zn), Calcium (Ca), Magnesium (Mg), Manganese (Mn), and Copper (Cu).

\subsection{Antioxidants and Phenolics}
Total antioxidant capacity was determined using the Oxygen Radical Absorbance Capacity (ORAC) assay, read on a BioTek Synergy HT microplate reader. Total phenolic content was measured via the Folin-Ciocalteu colorimetric assay, and anthocyanins were quantified using the pH differential method.

\subsection{Microbiological Plate Counts}
For food safety validation, microbial loads on leaves and growth pillows were quantified. Tissues were homogenized, serial diluted, and plated in duplicate. Aerobic Plate Counts (APC) were evaluated on Plate Count Agar incubated at 30$^\circ$C, and Yeast and Mold Counts (YMC) were evaluated on Potato Dextrose Agar incubated at 25$^\circ$C. Colony-forming units (CFU) were counted after 48-72 hours and converted to $\log_{10}$ CFU/g dry weight.

\section{Machine Learning Pipeline}

The analytical pipeline was constructed using Python 3.10, leveraging \texttt{scikit-learn}, \texttt{pandas}, and \texttt{shap}. 

\subsection{Dimensionality Reduction}
Principal Component Analysis (PCA) was employed for exploratory data analysis and unsupervised clustering. The elemental data matrix $X$ was standard-scaled to zero mean and unit variance. The covariance matrix $C$ was computed, and the eigenvalues $\lambda$ and eigenvectors $v$ were derived by solving:
\begin{equation}
C v = \lambda v
\end{equation}
where $C = \frac{1}{n-1} X^T X$. The first two principal components were extracted to visualize the variance captured by the multielemental profile.

\subsection{Classification and Validation Strategies}
To evaluate spaceflight signatures across crops, we trained classifiers to distinguish Flight vs. Ground Control samples based on elements and biochemicals. We implemented two validation designs:
\begin{itemize}
    \item \textbf{Leave-One-Mission-Out Cross-Validation (LOMOCV):} The model was iteratively trained on four spaceflight missions (e.g., VEG-01A, VEG-01B, VEG-03A, VEG-04A) and validated on the remaining mission (e.g., VEG-04B), preventing overfitting to mission-specific microenvironments.
    \item \textbf{Leave-One-Crop-Out Cross-Validation (LOCOCV):} To test if spaceflight physiological signatures are universal or crop-specific, the models were trained entirely on one species (e.g., Lettuce, OSD-745) and evaluated on the other (e.g., Mizuna, OSD-655).
\end{itemize}
We trained both Random Forest classifiers \cite{breiman2001} and TabPFN tabular foundation models \cite{hollmann2025} under these designs.

The Random Forest utilized an ensemble of decision trees, minimizing Gini impurity:
\begin{equation}
G = 1 - \sum_{i=1}^{c} p_i^2
\end{equation}
where $p_i$ is the probability of a sample belonging to class $i$ at a given node.

\subsection{SHAP Feature Importance}
To demystify the "black box" nature of the Random Forest, SHapley Additive exPlanations (SHAP) \cite{lundberg2017} were utilized. SHAP values calculate the marginal contribution of each elemental feature to the model's prediction across all possible coalitions of features:
\begin{equation}
\phi_i(f, x) = \sum_{S \subseteq N \setminus \{i\}} \frac{|S|!(|N| - |S| - 1)!}{|N|!} \left[ f(S \cup \{i\}) - f(S) \right]
\end{equation}
where $N$ is the set of all features, $S$ is a subset of features excluding feature $i$, and $f$ is the model prediction. This provided a ranked list of elements driving the FLT vs. GC distinction.

\subsection{Tabular Foundation Model Benchmarking}
To validate the Random Forest classification results, we implemented an independent benchmarking pipeline utilizing the state-of-the-art TabPFN (Tabular Prior-data Fitted Network) model \cite{hollmann2025}. Unlike traditional classification algorithms (e.g., Random Forests, Gradient Boosted Trees, or Multi-Layer Perceptrons) that require backpropagation and parameter tuning on the training dataset, TabPFN is a transformer-based foundation model trained offline using millions of synthetic datasets. 

At inference time, the model takes the entire training set (context) and the test set, performing a single forward pass (prior-data fitted learning) to generate zero-shot predictions. This approach is specifically optimized for small-sample classification tasks ($n < 10,000$ samples) and serves as a highly robust baseline to compare and contrast against the Random Forest classifier. Performance was evaluated using the identical LOMOCV strategy.

\subsection{Statistical Testing}
Non-parametric Mann-Whitney U tests were conducted for univariate comparisons. To account for multiple comparisons, $p$-values were adjusted using the Benjamini-Hochberg False Discovery Rate (FDR) procedure. Effect sizes were computed using Cohen's $d$:
\begin{equation}
d = \frac{\mu_1 - \mu_2}{s_{pooled}}
\end{equation}

\section{Interactive Dashboard Architecture}
To democratize access to the OSD-745 results, a static interactive dashboard was developed. The application was built utilizing HTML, CSS, and modern JavaScript, leveraging \texttt{Chart.js} and \texttt{Plotly.js} for responsive, client-side rendering of the PCA, SHAP plots, and boxplots. The dashboard incorporates the CoSE design system's aesthetic guidelines and is hosted via GitHub Pages for persistent, serverless accessibility.
